\documentclass[11pt]{article}
\usepackage[margin=1in]{geometry}
\usepackage{amsmath,amssymb,amsthm,booktabs}
\usepackage[hidelinks]{hyperref}
\usepackage{xurl}
\setlength{\emergencystretch}{3em}
\newtheorem{proposition}{Proposition}
\newtheorem{theorem}[proposition]{Theorem}
\newcommand{\C}{\operatorname{C}}
\title{A crosscap obstruction to twisted-link membership}
\author{Primary-source comparison and exact finite inputs}
\date{17 November 2026}
\begin{document}
\maketitle

The standard twisted-link theorem of Baader--Blair--Kjuchukova imposes
a crosscap restriction that excludes twelve retained sharp-four examples.
The restriction concerns every possible twisted diagram of a knot,
including diagrams with more crossings and arbitrary band signs.
It does not exclude rational-tangle extensions of the theorem or
arbitrary maximal-rank Coxeter quotients.

\section{A consequence of the bridge formula}

A twisted link in \cite{bbk} bounds a checkerboard surface $F$ built from
disks and twisted bands. Its weighted planar fat-vertex graph $\Gamma$
has one vertex per disk and one edge per band. All band weights have
absolute value at least two, and the plane dual $\Gamma^*$ is simple.
Theorem~1 of that paper states
\[
 b(K)=\#\{\text{regions of }\mathbb R^2\setminus\Gamma\}.
\]
Here $b$ is the actual bridge number, rather than the number of
maxima of a particular diagram.

\begin{proposition}\label{prop:crosscap}
If a knot $K$ is a twisted link in the precise sense of \cite{bbk}
and $b(K)$ is even, then
\[
                         \C(K)\le b(K)-1,
\]
where $\C(K)$ is the least nonorientable genus of a connected spanning
surface for $K$. In particular a twisted knot of bridge number four
has $\C(K)\le3$.
\end{proposition}
\begin{proof}
The connected surface $F$ retracts to the connected planar graph
$\Gamma$. Write $v,e,f$ for its vertex, edge and complementary-region
counts, including the unbounded region. Euler's formula and the
imported bridge formula give
\[
 \chi(F)=v-e=2-f=2-b(K).
\]
The surface has one boundary component because its boundary is the knot.
An orientable connected surface with one boundary component has
Euler characteristic $1-2g$, which is odd.
For even $b(K)$ the displayed Euler characteristic is even;
therefore $F$ is nonorientable.
Its nonorientable genus is
$1-\chi(F)=b(K)-1$, providing the required upper bound.
\end{proof}

This argument needs neither a classification of four-vertex Coxeter
graphs nor a bound on the crossing number of a twisted diagram.
It uses the topology of the spanning surface whose existence is part
of the standard twisted-link hypothesis.

\section{Twelve alternating examples}

Write a knot's Jones polynomial as
\[
 J_K(t)=\alpha_Kt^n+\beta_Kt^{n-1}
       +\cdots+\beta'_Kt^{s+1}+\alpha'_Kt^s,
 \qquad T_K=|\beta_K|+|\beta'_K|.
\]
Missing intermediate coefficients are understood to be zero.
Kalfagianni--Lee's Theorem~1.1 \cite{kl}
states, for a non-split prime alternating link of $k$ components
other than a $(2,p)$ torus link,
\[
                  \left\lceil\frac{T_K}{3}\right\rceil+2-k
                    \le \C(K).
\]
The rounding is a \emph{ceiling}. For knots the added term is one.

\begin{theorem}\label{thm:cases}
Each knot in the table below has the indicated Jones coefficient sum
and crosscap lower bound. Each has separately certified
$b(K)=\operatorname{mr}(K)=4$, and none is a twisted link in the
standard sense of \cite{bbk}.
\begin{center}
\small
\begin{tabular}{p{9.1cm}rr}
\toprule
knots & $T_K$ & lower bound on $\C(K)$\\
\midrule
$13a_{373},13a_{375},13a_{391},13a_{392}$ &10&5\\
$13a_{3264},13a_{3266},13a_{3275},13a_{3276}$ &9&4\\
$13a_{3145},13a_{3147},13a_{3148},13a_{3149}$ &12&5\\
\bottomrule
\end{tabular}
\end{center}
\end{theorem}
\begin{proof}
The finite input consists of the twelve literal thirteen-crossing
PDs in \path{twisted_link_crosscap_inputs.json}. Their named identities
are bound literally to the previously checked census/topology inputs.
The first eight bridge equalities use the seven-sheeted cover
certificates; the last four use the retained periodic-family seed
certificates at parameter zero. Those are separate dependencies.

Tracing the opposite half-edge through every crossing proves that
each PD has one component and alternating over/under encounters.
For each projection graph, exhaustive checking of all vertex
bipartitions, modulo complementation, finds minimum edge-cut size four.
If a Jordan curve meeting the diagram in two edges separated crossings
on both sides, those two edges would give a cut of size two.
Thus these connected alternating diagrams are prime.
The standard diagram of a $(2,13)$ torus knot has checkerboard vertex
counts $\{2,13\}$ and $T_K=1$; neither occurs here.
The passage from this prime alternating diagram criterion to the
non-split prime non-torus hypotheses is explicitly stated in
\cite[printed p.~2]{kl}, crediting Menasco. Its original proof is
imported at that scope.

For completeness, the coefficient computation is a full exact
Kauffman bracket state sum. At each crossing choose the pairings
$(a,b),(c,d)$ or $(b,c),(d,a)$ of the cyclically ordered PD ports.
For a state $s$, let $j(s)$ count choices of the second pairing and
let $\ell(s)$ be the number of resulting circles. The bracket is
\[
 \langle D\rangle=
 \sum_{s\in\{0,1\}^{13}}
 A^{13-2j(s)}(-A^2-A^{-2})^{\ell(s)-1}.
\]
Every one of the $8192$ states is evaluated with integer arithmetic.
Circle counts come from connected components of the smoothing.
The full Laurent polynomial is retained for each input.
Writhe normalization changes only an overall sign and monomial,
and substituting $t=A^{-4}$ therefore identifies the absolute
second and penultimate Jones coefficients with bracket coefficients
four exponent units from the two extremes.

For the first row both absolute coefficients are five.
For the second row they are four and five, in either order.
For the third row both are six.
The integer sums give the displayed $T_K$ values.
An independent reconstruction of both checkerboard graphs agrees:
if $v$ is the vertex count and $e'$ the edge count after parallel
edges are identified, the two values $e'-v+1$ match the two
computed absolute coefficients. The proof uses the full bracket
sum; this graph agreement is an additional consistency check.

Kalfagianni--Lee's inequality now supplies the displayed lower bounds,
all strictly greater than three. Proposition~\ref{prop:crosscap} would
give $\C(K)\le3$ if any of these bridge-four knots had a twisted diagram.
This contradiction proves the exclusions.
\end{proof}

\section{Source and verification scope}

The imported BBK statement was read in the published article's
Theorem~1 and its Section~4 proof, including the disk-band definition.
The full published HTML is retained as
\path{sources/bbk-coxeter-published-2021.html}.
The KL ceiling inequality and the following Menasco paragraph were
read in arXiv:1408.4493v3, printed pp.~1--2, and compared with the
published article's printed pp.~309--310 through the primary
published PDF. The valid cached source is
\path{sources/kalfagianni-lee-crosscaps-1408.4493v3.pdf}.
The full imported proofs have not been independently re-established.

The standard-library replay is
\path{verify_twisted_link_crosscap.py}.
It checks the literal PD bindings and the separately recorded
bridge-four conclusions; it does not rerun their earlier group
computations. It checks all diagram hypotheses and all complete
state sums, compares the full retained polynomials, and reconstructs
the checkerboard coefficients independently.
The receipt \path{twisted_link_crosscap_verification.json}
seals the actual inputs and source digests.
Known trefoil and figure-eight bracket polynomials are positive
controls. A changed bracket, a composite diagram falsely presented
as prime, and a changed source digest are rejected.

This is an internal exact finite replay with explicit classical
imports, not formal verification or independent specialist review.
Theorem~\ref{thm:cases} excludes all standard twisted diagrams for
these twelve knots; it is stronger than failure of a search among
minimal diagrams. The crosscap argument does not address the sixteen
nonalternating new seven-sheeted examples, all positive-parameter
members of the periodic families, or rational-tangle replacements.

For all eight alternating seven-sheet cases, the separate proof in
\path{UNIVERSAL_REFLECTION_OBSTRUCTION.tex} excludes every rank-four
Coxeter reflection quotient, including infinite targets. Hence none
admits a rank-four-preserving rational replacement retaining BBK's
reflection-label lower certificate. The corresponding reflection
comparison for the four earlier seeds remains separate.

The independent proof in \path{RATIONAL_REPLACEMENT_TOPOLOGY.tex}
now excludes all twenty-eight retained mixed-cover cases from arbitrary
rational replacements of twist boxes on any sharp-four BBK base.
That argument uses the complete exceptional-filling classification for
the octahedral double cover, rather than crosscap number.
The same note now gives a bridge-eight twisted base for $13a_{373}$
and proves that arbitrary rational replacements on unrestricted
higher-face bases are universal. Those geometric origins therefore
cannot be excluded; they do not preserve a sharp-four lower proof.
Nonsymmetric augmentation methods and worldwide priority remain
separate questions.

\begin{thebibliography}{9}
\bibitem{bbk}
S. Baader, R. Blair and A. Kjuchukova,
Coxeter groups and meridional rank of links,
Mathematische Annalen 379 (2021), 1533--1551.
Theorem~1 and Section~4.
\url{https://doi.org/10.1007/s00208-020-02124-z}.
\bibitem{kl}
E. Kalfagianni and C. Ruey Shan Lee,
Crosscap numbers and the Jones polynomial,
Advances in Mathematics 286 (2016), 308--337.
Theorem~1.1, printed p.~309; preceding/following diagram scope;
author preprint arXiv:1408.4493v3, printed pp.~1--2.
\url{https://doi.org/10.1016/j.aim.2015.09.017}.
\url{https://arxiv.org/abs/1408.4493v3}.
\end{thebibliography}
\end{document}
